\section{Skill vs. Luck: Can Past Winners Keep Winning?}\label{sec:skill}

\subsection{How Persistent Is Forecasting Skill?}

Wallets that perform exceptionally well in historical prediction markets may possess superior forecasting ability. However, favorable random outcomes can also produce impressive historical records, particularly when wallets participate in only a limited number of markets. To examine whether forecasting skill persists over time, this study compares the historical forecasting performance of individual wallets with their subsequent performance. Among 14,870 wallets that participate in at least five historical markets and five future markets, the correlation between historical and future forecasting scores is approximately 0.0234.

A permutation test is conducted to assess the statistical significance of this relationship. By controlling for differences in the number of forecasting opportunities available to individual wallets, the test constructs a reference distribution against which the observed correlation can be evaluated. The adjusted correlation is approximately 0.0227. Based on 1,000 permutations, the observed relationship is statistically significant, with a $p$ value of approximately 0.007. The results provide statistical evidence of persistent forecasting skill in prediction markets. Historical forecasting performance contains information about subsequent performance, indicating that a measurable component of forecasting ability persists across different market periods.

\studyfigure{figures/02_skill_persistence.png}{Skill persistence and permutation null}{Historical skill deciles and coverage adjusted permutation test for 14,870 wallets with at least five historical and five future markets. Observed adjusted correlation: 0.0227; permutation $p=0.007$ (seven day outcome availability specification).}{fig:skill_persistence}

\subsection{Experience Matters}

The persistence of forecasting skill may depend on the amount of historical experience available for evaluating each wallet. A longer forecasting record provides more observations from which individual performance can be estimated, potentially allowing genuine skill to be distinguished more clearly from random variation. To investigate this relationship, wallets are grouped according to the number of historical markets in which they have made valid predictions. The future forecasting performance of historically successful and unsuccessful wallets is then compared within each experience group.

The results reveal a more pronounced forecasting advantage among experienced wallets. For the 1,367 wallets with at least 50 historical markets and 50 future markets, the top 10\% based on historical forecasting performance achieve an average future forecasting score of approximately 0.0405. The corresponding score for the bottom 10\% is approximately 0.0025, yielding a difference of 0.0380. The permutation test further confirms the statistical significance of this difference, with a $p$ value of approximately 0.001. These findings highlight the importance of experience in identifying Smart Money. Wallets with more extensive forecasting records provide stronger evidence about their forecasting ability, making persistent skill more distinguishable in historical transaction data.

\studyfigure{figures/03_experience.png}{Forecasting skill by experience group}{Future market adjusted forecasting score for wallets in the top and bottom historical deciles, grouped by historical market coverage. Wallet counts for each group are shown along the horizontal axis.}{fig:experience}

\subsection{Can Empirical Bayes Identify Better Winners?}

The presence of persistent forecasting skill raises a practical question: how can wallets with superior forecasting ability be identified more accurately from historical observations? A conventional approach is to rank wallets according to their historical forecasting scores. However, extreme historical performance may reflect both genuine forecasting ability and favorable random outcomes. This phenomenon, commonly known as the Winner's Curse, is particularly relevant when individual wallets have limited forecasting histories. To improve wallet selection, this study employs Empirical Bayes estimation. The method adjusts observed forecasting scores according to the number of historical observations and their estimation uncertainty. Wallets with limited historical evidence receive greater adjustment toward the population average, while those with more extensive records retain a larger proportion of their individual performance estimates.

The results show that Empirical Bayes estimation substantially improves the identification of wallets with stronger future forecasting performance. Under the seven day outcome availability specification, the top 10\% of wallets selected by conventional historical ranking achieve an average future forecasting score of approximately 0.0130. The corresponding score increases to 0.0188 under Empirical Bayes ranking, representing an improvement of approximately \textbf{44\%}. Interestingly, only about half of the wallets selected by the two ranking methods overlap. The statistical adjustment therefore produces a substantially different group of potential Smart Money wallets. A similar improvement is observed under the 30 day specification, where Empirical Bayes ranking increases the average future forecasting score of the selected top 10\% by approximately 79.5\%.

\studyfigure{figures/04_empirical_bayes.png}{Raw and EB Top 10 percent future scores}{Future forecasting score of the selected top 10\% wallets using raw and Empirical Bayes rankings across minimum historical market counts. Seven day outcome availability specification.}{fig:eb}

These findings demonstrate that the identification of Smart Money depends not only on the presence of forecasting skill but also on how that skill is estimated. By accounting for the uncertainty associated with historical performance, Empirical Bayes provides a more effective approach to selecting wallets with stronger future forecasting outcomes.
